\section{A cyclic extension restriction obstruction}
\label{sec:rank11-c4}

\paragraph{Theorem.}
The based ring \texttt{R11\_1\_153} admits no categorification over
$\mathbb C$. The argument applies to every complex fusion category,
without a braiding, pivotality, or unitarity assumption.

\paragraph{A general restriction criterion.}
Let $\mathcal D$ be a fusion category, $G$ a finite group, $K\le G$, and
$c:G\to\operatorname{BrPic}(\mathcal D)$ a fixed homomorphism.
Put $A=\operatorname{Inv}Z(\mathcal D)$ with the action induced by $c$.
Suppose a reference $G$-extension realizing $c$ exists and
\[
 \operatorname{res}:H^2(G,A)\longrightarrow H^2(K,A)
 \quad\text{is zero.}
\]
Then the $K$-subextensions of all $G$-extensions realizing $c$ have
isomorphic based fusion rings, respecting their gradings and neutral
subcategories up to the chosen identifications.

Indeed, systems of multiplication functors form an $H^2(G,A)$-torsor
by \cite[Theorem~8.5]{ENOHomotopy}. Comparing two existing extensions,
their product systems differ by an $A$-valued $2$-cocycle. Restricting
the multiplication functors restricts this cocycle. If its restricted
class is zero, componentwise bimodule equivalences remove the
difference, so the restricted fusion coefficients agree. The remaining
$H^3(G,\mathbb C^*)$-freedom changes associators by scalars and leaves
multiplication functors unchanged
\cite[Theorem~8.9]{ENOHomotopy}. Duality is determined by the coefficient
of the unit, so it is preserved. This argument compares existing
extensions and requires no separate obstruction-vanishing assertion.

\paragraph{Exact structural hypotheses.}
Both \texttt{R11\_1\_104} and \texttt{R11\_1\_153} have dimension type
$(1^4,2^7)$, universal grading $C_4$, and homogeneous ranks $(5,2,2,2)$.
Their neutral ring is the Tambara--Yamagami ring on $V_4$; write $m$ for
its dimension-two simple. In a generating component there is a simple
$X$ with
\[
 XX^*=\mathbf1+h_L+m,\qquad X^*X=\mathbf1+h_R+m,
 \qquad h_L\ne h_R\ne\mathbf1,\quad h_L\ne\mathbf1.
\]
The source verifier derives the adjoint subring and all homogeneous
components directly from the tensors and checks every homogeneous
product, rather than only products of chosen representatives.

\paragraph{The neutral category is necessarily $\operatorname{Rep}(D_8)$.}
Every categorification of the neutral ring is
$\operatorname{TY}(V_4,\chi,\tau)$, with $\chi$ a nondegenerate symmetric
bicharacter and $\tau=\pm\tfrac12$.
The canonical internal-end algebras $XX^*$ and $X^*X$ have the form
$\mathbf1+h+m$. Their multiplication maps $h\otimes m\to m$ and
$m\otimes h\to m$ are nonzero: the $h$-summand arises from an invertible
stabilizing action on $X$. Associativity on $(h,m,h)$ therefore forces
$\chi(h,h)=1$. There are two distinct nonzero isotropic elements,
$h_L,h_R$. A nonalternating nondegenerate symmetric form on
$\mathbb F_2^2$ has exactly one, so $\chi$ is alternating.
The two remaining TY categories are $\operatorname{Rep}(D_8)$ and
$\operatorname{Rep}(Q_8)$, with groups of order eight.
But $\operatorname{BrPic}(\operatorname{Rep}(Q_8))\cong S_3$
\cite[Remark~9.4]{NikshychRiepel}. Any homomorphism from $C_4$ to this
group kills its degree-two element, forcing that homogeneous bimodule
to be regular of rank five, contrary to its rank two.
Thus $\mathcal D\simeq\operatorname{Rep}(D_8)$. This includes all neutral
associators and both TY signs.

\paragraph{The Brauer--Picard homomorphism is unique up to neutral equivalence.}
By \cite[Section~9.1]{NikshychRiepel},
$\operatorname{BrPic}(\mathcal D)\cong S_4$.
Reindexing the alternating TY data realizes the six isometries of $V_4$
as tensor autoequivalences. Their map into the Brauer--Picard group is
injective: its kernel consists of inner tensor autoequivalences, which
are tensor isomorphic to the identity in a braided category.
This $S_3$ is the stabilizer of the regular left-module class, a point
$p$ of the four-point $S_4$ action.
All three nonidentity components have rank two, so the first, second,
and third powers of the grading generator move $p$. Its image is
therefore a $4$-cycle. The $S_3$ fixing $p$ conjugates the six $4$-cycles
transitively: its intersection with the centralizer of one is trivial.
Consequently every relevant $c$ can be aligned with the reference by
an actual tensor autoequivalence of $\mathcal D$.

\paragraph{The coefficient module.}
Here $A=\operatorname{Inv}Z(\mathcal D)\cong C_2^3$ and, as an $S_4$-module,
\[
 A\cong\mathbb F_2^4/\langle(1,1,1,1)\rangle.
\]
For completeness this action can be identified intrinsically. Let $z$
be the nontrivial central invertible with underlying tensor unit: its
half-braiding acts as $+1$ on the four invertibles of $\mathcal D$ and
$-1$ on $m$. Every tensor autoequivalence fixes $z$, so its $S_4$ orbit
has size one or four. If its orbit had size one, the $C_4$ action of
the explicit \texttt{R11\_1\_104} model would fix $z$.
Relative-center equivariantization \cite{CentersGraded} would then lift
$z$ to a central invertible with underlying unit: the projective
linearization obstruction lies in $H^2(C_4,\mathbb C^*)=0$.
But tensor automorphisms of the identity of that category are
characters of its universal grading $C_4$, hence restrict trivially
to its adjoint category $\mathcal D$. This contradicts the nontrivial
restriction $z$.
Thus the orbit consists of four distinct nonzero vectors. They span
$A$, since a two-dimensional space over $\mathbb F_2$ has only three
nonzero vectors. The unique relation is $S_4$-invariant and is their
sum zero, proving the displayed module identification.

\paragraph{Vanishing of the restriction.}
Let $t$ be a $4$-cycle and $N=t-1$ on $A$. This is a single unipotent
Jordan block of size three:
\[
 N^3=0,\qquad \operatorname{rank}N=2,\qquad
 \operatorname{rank}N^2=1,\qquad \ker N=\operatorname{im}N^2.
\]
The $C_4$ norm is $1+t+t^2+t^3=N^3=0$, so cyclic cohomology gives
\[
 H^2(C_4,A)=\ker N,\qquad
 H^2(\langle t^2\rangle,A)=\ker N^2/\operatorname{im}N^2.
\]
Even-degree restriction sends an invariant representative to itself.
Since $\ker N=\operatorname{im}N^2$, this restriction is zero.
This computation is also verified without invoking the periodic
resolution: an independent normalized bar-cochain program enumerates
all $128$ normalized $C_4$ $2$-cocycles, finds the $64$ coboundaries,
and supplies a $C_2$ $1$-cochain trivializing the restriction of every
cocycle. It also verifies the six conjugacy representatives above.

\paragraph{Reference model and contradiction.}
The reference \texttt{R11\_1\_104} is the explicit group-theoretical
category $\mathcal C(G,\omega,H,1)$ provided by the finite certificate
\texttt{jordan\_gt\_model.json}. Here $|G|=32$, $H\cong V_4$, and
$\omega$ is a normalized $\mu_8$-valued $3$-cocycle trivial on $H$.
The group is the nonsplit Jordan extension of $C_2^3$ by $C_4$,
with $t^4$ equal to the fixed nonzero vector.
The certificate includes the complete group table and cocycle,
all simple labels, the full model fusion tensor, and a unit-preserving
basis bijection to the original census tensor. All $32^4$ cocycle
identities and all fusion coefficients are checked. An independent
implementation reconstructs all $1331$ coefficients from free algebra
bimodules, balanced tensor quotients, and graded intertwiner systems;
it imports no projective-character convolution routine. This audit
passes and is executed by the exclusion verifier.

In \texttt{R11\_1\_104}, the degree-two simples have source labels
$9,10$ and are both self-dual. In \texttt{R11\_1\_153}, those labels
are $1,2$ and they are dual to each other. Thus the respective even
subrings have seven and five self-dual simples and are not isomorphic
as based rings. A hypothetical categorification of the latter ring
has the forced neutral category and the same $4$-cycle homomorphism
as the reference, up to neutral equivalence. The restriction criterion
forces its even subring to equal that of the reference, a contradiction.

\paragraph{Scope and reproducibility.}
The exclusion uses the new general restriction criterion and a positive
reference model, not a claim that a bounded model search is exhaustive.
The earlier failed cocycle ans\"atze have no role in the proof.
The source-bound records are \texttt{c4\_reduction.json} and
\texttt{c4\_exclusion.json}; the independent normalized-cochain audit is
\texttt{c4\_torsor\_independent.json}. The categorical neutral-module
audit is \texttt{C4\_NEUTRAL\_MODULE\_AUDIT.md}.
